% STEADY_STATE_ADDENDUM.tex, the scope of the main paper with respect to the long run: what its
% one-shot Nash analysis covers, and what a contest that repeats, with the incumbent and the
% challengers exchanging roles, would decide. Not a literature position (author, 15 October 2026:
% "the goal of the addendum is to just establish the scope of the main paper"). Origin (author, 15 October 2026: "if we talk
% about long-term steady-state equilibria - then we have to consider the possibility that the
% incumbent becomes the contender and vice versa over time").
%
% RULES
%   A statement that rests on a row of skeleton_who_pays.tex cites the row. A statement that follows
%   from a row in one step says so and names the step. Everything else is a conjecture or an
%   open question, labelled S-n, with the world in which it is false. Terms are the manuscript's;
%   checks/verify_docs.py (D5) checks the terms and that every cited row exists.

\documentclass[11pt]{article}
\usepackage[T1]{fontenc}
\usepackage[utf8]{inputenc}
\usepackage{mathptmx}
\usepackage{microtype}
\usepackage[a4paper,margin=2.2cm]{geometry}
\usepackage{amsmath,amssymb}
\usepackage{booktabs}
\usepackage{array}
\usepackage{longtable}
\usepackage[round]{natbib}
\usepackage{parskip}
\usepackage[hidelinks]{hyperref}

\newcolumntype{P}[1]{>{\raggedright\arraybackslash}p{#1}}

\title{The Long Run: Scope of the Main Paper}
\author{Companion note to the manuscript}
\date{15 October 2026}

\begin{document}
\maketitle

\section{The scope}

The paper is a Nash-equilibrium analysis. Every equilibrium in it is an
equilibrium of a game played once: the incumbent's identity is fixed, the prize
is awarded, and the model stops (manuscript, Appendix~D, ``The contest''). Over
time the contest repeats. The winner holds the position next period, the
displaced holder competes again, and wealth responds to past success.

The two are read together as forces and states. At any date the state is the
wealth profile, who holds the position, and the prize, the fee and the weight in
force. The one-shot equilibrium at that state is where the opposing forces of
the model settle. The fee weighs less on the rich, so they enter first (K1,
K3); the gain from entering falls as more enter, which caps entry (M4, K4); a
single-crossing widening of wealth moves the count in a direction set by where
the cut-off wealth stands against its break-even wealth (K2); a heavier weight on the base component lowers entry (K8)
without guaranteeing a better winner (M32). The long run decides which states
are reached. Its variables are not free parameters but values the process
arrives at, and every claim of the paper holds at whatever values it reaches.
The claims are therefore conditional on the state, which is the scope this note
records.

\paragraph{The bridge condition.} The reading holds when the period game is the
one-shot game. That is so if agents are myopic, or if the value of holding the
position after this period is the same for every player. A forward-looking
challenger values winning at $V$ plus the discounted value of incumbency. Since
$V$ enters the gain only as a scale (M1), a continuation value common to all
players changes the prize and nothing else, and every result holds with the
larger prize. If the continuation value depends on a player's own wealth, the
prize is no longer valued alike by rich and poor (Appendix~D, ``Payoffs''), and
wealth acts through the prize as well as through the fee. The period game is
then a different game, and that is the one case in which the long run would
change the claims themselves rather than the state at which they are read.

\section{What the one-shot model fixes}

Given the wealth profile and the incumbent at a date, the one-shot results hold
at that date. The table lists the ones the extension would lean on.

\begin{center}
\small
\begin{tabular}{@{}P{4.4cm}P{8.2cm}P{2.6cm}@{}}
\toprule
\textbf{Statement} & \textbf{Content} & \textbf{Rests on} \\
\midrule
The count each period & Given the profile, every one-shot equilibrium has
  $k^*$ entrants, and the $k^*$ richest entering is one of them. & K1, M6, M8 \\
  \addlinespace
Wealth enters only through $\kappa$ & The fee's utility cost
  $\kappa(w)=u(w)-u(w-c)$ falls with wealth. & K3, M23 \\ \addlinespace
The prize is outside $u$ & The prize enters payoffs linearly and outside the
  utility of wealth, so it does not change wealth. & Appendix~D, Payoffs \\
  \addlinespace
The incumbent's entry & As player $0$ with the lowest cost, she is in every
  non-empty equilibrium outside the band of M6. & M30 \\ \addlinespace
The entrants' share & The prize goes to an entrant with probability $\xi(k)$,
  non-decreasing in $k$, equal to $1$ when all enter, non-increasing in $\mu$. &
  M31 \\ \addlinespace
A common bonus & It washes out when every player holds it and favours holders
  once some stop buying. & M32 \\
\bottomrule
\end{tabular}
\end{center}

\section{The ingredients of a steady state}

Three choices turn the one-shot game into a process. Each is a modelling fork,
and the extension must take one branch of each.

\paragraph{Turnover.} The winner at date $t$ is the incumbent at $t+1$; the
displaced holder joins the challengers. Variants: the position lasts a fixed
number of periods, or a replacement shock ends tenure.

\paragraph{The law of motion of wealth.} In the manuscript the prize is outside
$u$, so winning changes nothing but the payoff. For wealth to respond to
success the position must pay. Two branches.
\begin{itemize}
\item \textbf{S-A, the position pays.} The holder receives an income $y>0$ each
  period of tenure, in wealth units, besides the prize's value $V$ in utility.
  Then $w_{t+1}=w_t-c\,e_t+y\,\mathbf 1\{\text{holds}\}$, plus any endowment
  shock. This is a step from the privilege contest of Appendix~D towards a
  rent-seeking contest, since part of what is won is now money.
\item \textbf{S-B, the position does not pay.} Wealth moves only through the
  fees paid and an exogenous endowment process. Then every entrant is poorer
  next period by $c$, and success does not feed wealth.
\end{itemize}

\paragraph{The life of the components.} Either the base and bought components
are redrawn each period (memoryless), or the bought component persists for the
holder once bought (a credential kept). Under persistence the fee is paid once,
which is an investment rather than a fee, and the incumbent is scored as an
entrant by inheritance rather than by choice.

\section{How the claims change}

The one-shot claims survive as statements about each period, since each period
is the one-shot game at the current state. What changes is the state. The
table gives, for each headline, what is unchanged and what the steady state
adds, with the conjectures labelled.

\footnotesize
\begin{longtable}{@{}P{1.6cm}P{5.2cm}P{8.4cm}@{}}
\caption{Each headline, per period and in the steady state.}\\
\toprule
\textbf{Row} & \textbf{Per period} & \textbf{In the steady state} \\
\midrule
\endfirsthead
\toprule
\textbf{Row} & \textbf{Per period} & \textbf{In the steady state} \\
\midrule
\endhead
\bottomrule
\endfoot
K1 & Unchanged: the count is pinned at each date by the current profile. & The
  incumbent's identity becomes a Markov chain on the players. \textbf{S-1}
  (one step from M31 and the symmetry of the entrant law): with components
  redrawn each period and the incumbent scored as an entrant, the $k+1$
  entrants win with equal probability, so the incumbent is displaced with
  probability at least $1-\xi(k)/(k+1)$, which is at least $\tfrac12$ whenever
  one challenger enters, and with probability $1-\xi(0)$ when none does.
  Closure under redraw therefore needs $k^*=0$ and $\xi(0)=1$, which holds
  when $(1-\mu)s$ exceeds $\mu$ times the width of the support of $r$ almost
  surely, as at the $\mu=0$ boundary with $s>0$.
  \emph{False if} the entrant law is not exchangeable across entrants, which
  the anonymity of M26 rules out. \\ \addlinespace
K2 & Unchanged for a change of the current profile. & The change of profile is
  no longer exogenous. Under S-A the holder's wealth rises by $y$ each period, so the
  stationary profile is more dispersed than the endowment profile.
  \textbf{S-2} (conjecture): the stationary dispersion of wealth rises in $y$.
  \emph{False if} turnover is so fast that no holder accumulates, which S-1
  makes possible at small $\xi(k)$. The sign rule of K2 then applies to a
  single-crossing widening the process generates, read at the stationary
  cut-off wealth. \\
  \addlinespace
K3 & Unchanged. & Unchanged; $\kappa$ is still the only channel within a
  period. Under S-A wealth also records past success, so the richest-first
  sorting of K1 reads as a sorting by past success. \\ \addlinespace
K4, K5 & Unchanged. & The prize gains a second channel. \textbf{S-3}
  (conjecture, S-A): the stationary count is non-decreasing in $y$, because a
  displaced holder returns richer by her tenure's income, and by M5 a richer
  challenger enters whenever a poorer one does. \emph{False if} the income $y$ raises the holder's
  wealth enough to change the law of $s$ or $r$, which the anonymity of M26
  excludes, or if $y$ is taxed into the fee. \\ \addlinespace
K6 & Unchanged. & Unchanged; the number of challengers is fixed. \\
  \addlinespace
K7, K8 & Unchanged. & \textbf{S-4} (conjecture, persistence): with the bought
  component kept by the holder, a higher $\mu$ lowers the number who buy (K8),
  and the holder's kept bonus then counts against fewer holders, so tenure
  lengthens and mobility falls; closure can then arise at $\mu>0$. Under
  redraw there is no such persistence, and S-1 bounds mobility below.
  \emph{False if} the kept bonus depreciates faster than the entry response. \\
  \addlinespace
M30 & The incumbent chooses as player $0$. & \textbf{S-5} (one step from M5
  and the law of motion): under S-A each period of tenure adds $y$ to the
  holder's wealth while no challenger's wealth rises, so after enough periods
  of tenure the holder is the richest player, has the lowest cost and, by M5,
  enters whenever any challenger does. ``Scored as an entrant'' is then the
  outcome of the process for a holder of long enough tenure, not an
  assumption. \emph{False if} tenure ends before that, which S-1 makes likely
  under redraw whenever a challenger enters. \\ \addlinespace
M32 & Unchanged. & The common-bonus mechanism compounds. Under persistence, a
  holder who keeps the bonus wins against outsiders with better base
  components in every period in which some challengers have stopped buying, so
  the winner's base component in the steady state is below its one-shot
  level. \textbf{S-6} (conjecture): under persistence and S-A, the stationary
  expected base component of the holder falls in $\mu$. \emph{False if} S-1's
  turnover dominates, as under redraw. \\ \addlinespace
C1 to C14 & Unchanged; they are claims about the one-shot model. & Unchanged as
  claims about each period; any claim about the steady state itself would be a
  new claim. \\
\end{longtable}
\normalsize

\section{Mobility and closure as objects}

Mobility is the stationary probability that the position changes hands.
Closure is mobility zero. In the one-shot model closure is the $\mu=0$
boundary, where outsiders score nothing (manuscript, Research question). In the
process it has two further routes, each a conjecture above: a holder's wealth
that makes her the lowest-cost player every period (S-5, which gives
persistence of entry, not of the position), and a kept bonus that counts
against fewer holders as $\mu$ rises (S-4, which gives persistence of the
position). S-1 says that under redraw neither route closes the position while
any challenger enters. So the sharpest statement the extension can aim at is
that closure needs either the $\mu=0$ boundary or a bought component that
persists; wealth alone, under redraw, does not close it.

\section{What this gives the manuscript}

Two sentences in Appendix~D of the manuscript (``The contest''). The first says
that ``equilibrium'' means the one-shot equilibrium with the incumbent's
identity fixed, so the claims are conditional on the state, and that a steady
state is outside the model. The second gives the bridge condition, myopia or a
continuation value common to all players, under which each period of a
repeated contest is the one-shot game.

\bibliographystyle{plainnat}
\bibliography{refs}

\end{document}
